\subsection{Contexte de recherche et collaboration}
Cette étude s'inscrit dans le cadre d'un projet d'Introduction à la Recherche au sein de l'IMAG,
en collaboration avec les équipes \textbf{DANCE}\footnote{\url{https://team.inria.fr/dance/}}
avec Féderica Garin\footnote{\url{https://www.gipsa-lab.grenoble-inp.fr/~federica.garin/}} et
Paolo Frasca\footnote{\url{https://www.gipsa-lab.grenoble-inp.fr/~paolo.frasca/}} et
\textbf{POLARIS}\footnote{\url{https://team.inria.fr/polaris/}} avec Nicolas Gast\footnote{\url{https://polaris.imag.fr/nicolas.gast/}}
de l'INRIA Grenoble.
L'équipe POLARIS travaille sur l'analyse des performances des très grands systèmes distribués,
tandis que l'équipe DANCE se concentre sur l'étude des dynamiques sur les réseaux.
Cette collaboration interdisciplinaire nous a permis d'aborder la question du risque systémique avec des outils et perspectives complémentaires,
en associant modélisation mathématique et simulation informatique.

\subsection{Le risque systémique et son importance}
L'existence des marchés financiers au cours de l'histoire du capitalisme se caracterise par les crises périodiques qui sont censés être des évenements extrême,
pourtant, on en voit plus souvent que prevu. Dans la plupart de nos modèles, on sous estime les interdépendances bancaires. Parmi les crises récente, l'envergure de la crise 2007-2008, le rôle joué par les interconnexions des banques dans cette crise sera un point de depart pour notre
problème.
Cette crise a fourni une illustration dramatique de la propagation des défaillances
dans un réseau financier \cite{peyton2016contagion}. Lehman Brothers était la 4e plus grande banque
d'investissement en Amérique. En se basant sur la description de \cite{jglass2006topologyus} du réseau
interbancaire aux USA, qui nous indique qu'il existe des grandes banques nommées \textbf{core},
très fortement liées entre elles, et les \textbf{périphéries} qui empruntent à ces grandes banques.
Ces banques avec un capital très important portent l'appellation "\textbf{too big to fail}".
Lehman Brothers faisait parti des \textbf{core} et même des core les plus important.
Ainsi lorsque Lehman Brothers fut en faillite, elle provoqua une cascade de défaut. Cependant,
d'autres facteurs économiques étaient présents, comme la liquidité et les produits dérivés.
La forte connectivité du réseau de dette interbancaire aux USA a créé des canaux de contagion qui ont
aggravé la situation de la crise. Ces événements ont révélé comment des institutions fortement
interconnectées pouvaient transformer des chocs initialement localisés en une crise systémique.


L'étude du risque systémique est devenue centrale dans l'analyse de la stabilité financière.
Cette crise a révélé comment la défaillance d'institutions individuelles pouvait se propager à
travers un réseau d'interconnexions et menacer l'intégrité du système financier global.
Cette crise systémique est justement ce qui va nous permettre de donner un sens au \textbf{risque systémique},
qui n'est rien d'autre que la probabilité qu'une crise systémique se produise.
Ainsi, pour un réseau de dette interbancaire donné,
si on réussit à trouver les paramètres qui influence le plus la fragilité du réseau,
on réussira à donner une mesure de ce risque.


\subsection{Modélisation et approche méthodologique}
Pour donner un sens à ces paramètres, parlons rapidement du modèle choisi.
Notre approche se situe à l'intersection de la théorie des graphes aléatoires,
de l'informatique appliquée à la finance et à l'économie.
On considère un réseau de dette interbancaire comme un graphe pondéré et orienté,
avec le sens des flèches représentant l'obligation de l'emprunteur vers le prêteur.
Supposons dans cette introduction que ce graphe représente le réseau interbancaire d'un monde économique fictif,
où les banques sont des entités qui n'échangent que des dettes,
il n'y a aucun problème de liquidité ou encore de taux d'intérêt.
Ainsi, on représente une banque avec un bilan simplifié, les dettes extérieures au réseau
et les dettes liées au réseau.


Un premier problème qu'on peut soulever, comme \cite{peyton2016contagion} le souligne dans son article,
pourquoi les banques décident de tant se lier à d'autres banques à travers les dettes ?
Ce que la théorie financière nous enseigne, c'est que toutes les banques cherchent à réduire
le risque derrière chaque opération financière c'est le risque global. Ainsi, la forte connectivité permettant de diversifier
leurs actifs chez différents acteurs, les banques empruntent et prêtent le plus possible.
Mais on verra par la suite que selon la nature du réseau et selon la longueur des cycles présents dans le réseau,
ces diversifications peuvent se retourner contre les banques en cas de défaut d'un de leurs créanciers.
On voit donc que le niveau d'interconnexion joue un rôle important dans la fragilité du réseau.


\subsection{Objectifs et problématique}
Dans cette étude, nous examinons spécifiquement l'impact du niveau d'interconnexion sur la résilience d'un réseau
financier. La fragilité d'un système financier est un sujet important tant d'un point de vue économique que financier.
Économiquement car plus le risque systémique est élevé, plus les conséquences sur l'économie réelle seront importantes,
se traduisant par des récessions, du chômage et des interventions publiques coûteuses.
Financierèment car dans les mécanismes internes du système, cela se traduit par une plus grande volatilité
et une diminution du volume de capital bancaire.


\subsection{État de l'art et positionnement}
L'impact de l'interconnexion sur la résilience d'un réseau est une rélation étudier par beaucoup de chercheur,
comme par exemple \cite{peyton2016contagion} dans son article Systemic risk in
financial systems, ils étudient l'importance non pas d'un seul paramètre,
mais plutôt comment plusieurs paramètres interagissent entre eux. En l'occurrence,
ils regardent comment les expositions de dettes ou l'effet de levier influence le risque systemique,
et montre l'ambivalence de l'interconnexion interbancaire,
qui peut à la fois diversifier les expositions aux risques mais aussi créer des canaux
par lesquels les chocs peuvent se propager par contagion. Une particularité de ce papier est qu'il synthétise l'état de l'art sur le risque systemique, et donne les différents approchent pour
modeliser et tester la résilience d'un réseau, quel paramètre serait le plus significatif. Ce qui nous a beaucoup inspiré
car notre approche consiste à explorer sa synthèse sur les paramètres, et montrer grace à nos simulation une ilustration empirique.

Dans la même approche, on a le papier de \cite{nier2008stability} qui fait une simulation
à travers un graphe fixé d'Erdos-Renyi, ce qui nous a beaucoup inspiré pour l'ouverture de ce texte.
Et pour finir sur les différentes approches qui existent, on peut aussi stipuler l'importance
que joue l'algorithme \cite{eisenberg2001systemic}, qui permet pendant une simulation,
ou n'importe quel cas où un réseau interbancaire intervient, de trouver un vecteur d'équilibre de paiement,
qu'on prendra le temps de définir dans la suite.


Dans ce contexte, plusieurs questions fondamentales émergent :
comment mesurer précisément le risque systémique ?
Quel rôle joue le volume du capital échangé entre les banques à travers les dettes dans ce risque systémique ?
Dans le cas d'un non-remboursement provenant de l'extérieur du système (choc exogène),
comment se propagent les défaillances en cascade à travers le réseau ?


\subsection{Problématique centrale et contribution}
C'est là que notre problématique centrale intervient :

dans quelle mesure les interconnexions permettent-elles de réduire l'impact d'un choc ?
À partir de quel seuil d'interconnexions ce choc cesse-t-il d'être dilué pour, au contraire, être amplifié ?
Parmi les multiples paramètres pouvant influencer la résilience d'un réseau financier
(volume du capital, structure hiérarchique, nature des actifs, etc.),
nous avons choisi de nous concentrer spécifiquement sur le degré d'interconnexion.
Ce choix se justifie par le paradoxe apparent qu'il soulève : alors qu'une plus grande diversification
des connexions devrait théoriquement diluer les risques,
l'expérience de la crise suggère qu'au-delà d'un certain seuil, elle pourrait au contraire les amplifier.
Cette question demeure centrale dans la littérature récente sur la contagion financière.


\subsection{Plan de l'article}
Cet article s'organise comme suit.
Nous formalisons le problème~\ref{sec:formalisation}, en présentant le modèle et puis on présente notre approche par simulation et annonce les résultats de la simulation.
Et enfin on analyse les résultats ~\ref{sec:analyse} par des analyses graphiques.
Nous discutons des implications de ces résultats et proposons
des pistes d'amélioration, notamment par l'intégration d'autres caractéristiques structurelles
comme l'intensité des obligations, le volume du capital et la concentration des nœuds.



